\chapter{Trigonometric Equations}
\label{ch:ch11-equations}

\section{From Identities to Equations}

The previous three chapters were concerned primarily with identities: how
rotations combine, how the unit circle constrains admissible trigonometric
values, and how apparently complicated formulas can be derived from a
small collection of fundamental relationships. Throughout that development
the task was always the same in kind --- start from one expression,
transform it into another, and show that the two are equivalent. The task
of this chapter is different. Rather than proving that two expressions are
always equal, this chapter asks for which values of an angle a particular
equation becomes true, moving from proving identities to solving equations.

At first glance this may look like a familiar algebraic exercise: solving
$x^2 - 5x + 6 = 0$ means finding the values of $x$ that make the equation
true, so why should a trigonometric equation be any different? The answer
is periodicity. A polynomial equation ordinarily has only finitely many
solutions, but trigonometric functions repeat indefinitely, and if a point
on the unit circle returns to the same position every full revolution,
then any trigonometric value that occurs once must occur again and again.
Trigonometric equations therefore typically possess infinitely many
solutions, and the challenge shifts from finding a single value that works
to describing the entire family of values that work. When solving an
algebraic equation it is often enough to list the roots; when solving a
trigonometric equation, the final answer is usually not a number but a
description of a repeating structure.

\section{The Simplest Equations}

The natural place to begin is with an equation involving a single
trigonometric function, such as $\sin\theta = a$. Chapter~7 introduced the
inverse sine function and showed how to find an angle whose sine equals a
specified value, giving $\theta = \arcsin(a)$; this angle is called the
\emph{principal value}, the particular solution returned by the inverse
function. It is rarely the complete answer. A horizontal line drawn across
the unit circle generally intersects it at two points sharing the same
vertical coordinate and therefore the same sine value, so
$\sin(\pi/6) = \tfrac12$ but also $\sin(5\pi/6) = \tfrac12$; the inverse
function returns only the first of these, while the geometry of the circle
guarantees that a second one exists. The periodicity of sine then produces
infinitely many further solutions, since adding any integer multiple of
$2\pi$ leaves the sine value unchanged.

\begin{theorem}
Let $\alpha = \arcsin(a)$. Every solution of $\sin\theta = a$ is given by
$\theta = \alpha + 2\pi k$ or $\theta = \pi - \alpha + 2\pi k$, for $k$ any
integer.
\end{theorem}

This theorem is often memorized as a formula, but its content is simpler
than it looks: the angle $\alpha$ names one point on the unit circle, and
$\pi - \alpha$ names its mirror image across the vertical axis, a point
with the same height and hence the same sine value; adding multiples of
$2\pi$ simply repeats both positions on every subsequent revolution. The
theorem formalizes the geometry of the circle rather than introducing
anything beyond it.

Cosine behaves similarly, though its symmetry runs the other way: equal
cosine values arise from reflection across the horizontal axis rather than
the vertical one. If $\alpha = \arccos(a)$, every solution of $\cos\theta =
a$ can be written as $\theta = \pm\alpha + 2\pi k$. Tangent differs because
its period is $\pi$ rather than $2\pi$, established in Chapter~6, so once a
single solution $\arctan(a)$ is known, every other solution is obtained by
adding integer multiples of $\pi$ alone: $\theta = \arctan(a) + \pi k$.

\begin{example}
Solve $\sin\theta = \tfrac12$ completely. The inverse sine function gives
the principal solution $\theta = \arcsin(1/2) = \pi/6$, but this alone
cannot be the whole story: the graph of sine shows that the value
$\tfrac12$ occurs twice within a single cycle, at the reference angle and
again at its Quadrant~II partner, $\theta = 5\pi/6$. Having found both
solutions in one revolution, periodicity extends each into an infinite
family:
\[
\theta = \frac{\pi}{6} + 2\pi k \qquad \text{or} \qquad \theta = \frac{5\pi}{6} + 2\pi k,
\]
for $k$ any integer. The pattern is worth naming explicitly, since nearly
every equation in this chapter follows it: begin with an inverse function,
use geometry to find any remaining solutions within one revolution, and use
periodicity to generate the entire family from those.
\end{example}

\section{A General Strategy}

The equations solved so far shared a convenient feature: the trigonometric
function was already isolated. Many equations do not arrive in so
convenient a form. The equation $2\sin\theta\cos\theta - \sin\theta = 0$,
for instance, offers no immediate use for an inverse trigonometric
function, since neither sine nor cosine has been isolated, and some
algebraic work must be done before any geometric reasoning can begin. This
observation is worth stating as a general principle: solving a
trigonometric equation is a collaboration between algebra and geometry,
where algebra simplifies the equation until its structure becomes visible
and geometry then identifies the corresponding angles and their periodic
repetitions. Neither approach is sufficient alone. The identity toolkit of
Chapters~8 through 10 exists largely for this purpose --- not merely as a
set of facts to be proved, but as a set of tools for rewriting an equation
into a form whose solutions can be read off geometrically.

\begin{algorithmproc}{How to Solve a Trigonometric Equation}
Begin by treating the equation exactly as an algebraic equation: factor
where possible, move all terms to one side, and look for common factors or
useful substitutions. If the equation mixes more than one trigonometric
function or more than one angle, apply an identity from Chapters~8 through
10 to reduce it to a single function of a single angle. Once the equation
involves only one trigonometric function, isolate that function and apply
the appropriate inverse function together with the reference-angle and
quadrant reasoning of Chapter~4 to find every solution within one
revolution. Extend those solutions to the general solution using
periodicity, then check any restricted domain and verify that no
extraneous solutions were introduced by the algebraic steps along the way.
\end{algorithmproc}

\begin{example}
Solve $2\sin\theta\cos\theta - \sin\theta = 0$. Factoring out the common
factor $\sin\theta$ gives $\sin\theta(2\cos\theta - 1) = 0$, and since a
product equals zero only when at least one factor does, either
$\sin\theta = 0$ or $2\cos\theta - 1 = 0$. The original problem has become
two simpler ones. The first, $\sin\theta = 0$, holds whenever the
corresponding point lies on the horizontal axis of the unit circle, giving
$\theta = k\pi$. The second, $\cos\theta = 1/2$, has reference angle
$\pi/3$; since cosine is positive in Quadrants~I and IV,
$\theta = \pm\pi/3 + 2\pi k$. Combining the two families gives the
complete solution set. Notice how little trigonometry was needed once the
factorization was found: the difficult step was algebraic, and the
geometry that followed was routine.
\end{example}

\section{Quadratic Trigonometric Equations}

Some trigonometric equations resemble polynomial equations even more
closely. The equation $2\sin^2\theta - \sin\theta - 1 = 0$ involves only
powers of a single trigonometric function, and looks remarkably like
$2x^2 - x - 1 = 0$; this resemblance is not accidental. Whenever an
equation involves only one trigonometric function and its powers,
substitution converts the problem into an ordinary algebraic equation: set
that quantity temporarily equal to a new variable, solve the resulting
polynomial equation, and translate the roots back into trigonometric
language once the algebra is finished.

\begin{example}
Solve $2\sin^2\theta - \sin\theta - 1 = 0$. Substituting $u = \sin\theta$
gives the quadratic $2u^2 - u - 1 = 0$, which factors as
$(2u+1)(u-1) = 0$, so $u = 1$ or $u = -\tfrac12$. Substituting back,
$\sin\theta = 1$ or $\sin\theta = -\tfrac12$. The equation $\sin\theta = 1$
corresponds to the top of the unit circle, giving $\theta = \pi/2 + 2\pi
k$. The equation $\sin\theta = -\tfrac12$ has reference angle $\pi/6$; since
the value is negative, the relevant quadrants are III and IV, giving
$\theta = 7\pi/6 + 2\pi k$ and $\theta = 11\pi/6 + 2\pi k$. The complete
solution consists of these three infinite families together. What matters
about this example is not the particular answer but the pattern: many
trigonometric equations are solved by temporarily converting them into
algebraic equations, solving those, and translating the roots back into
geometric language.
\end{example}

\section{Why the Identity Toolkit Was Necessary}

A natural question arises at this point: why spend three chapters
developing identities if factoring and substitution already solve so much?
The answer is that many equations do not initially involve a single
trigonometric function. The equation $\cos(2\theta) + \sin\theta = 0$
contains two different trigonometric expressions, and neither factoring
nor a direct inverse function is immediately useful against it; the
obstacle is not the equation itself but its form, and Chapter~10 supplies
exactly the tool needed to change that form.

\begin{practicenow}
Before reading the worked example below, try rewriting $\cos(2\theta) +
\sin\theta = 0$ using the double-angle identity $\cos(2\theta) = 1 -
2\sin^2\theta$, and see how far the substitution $u = \sin\theta$ takes
you toward a solution.
\end{practicenow}

\begin{example}
Solve $\cos(2\theta) + \sin\theta = 0$. Using
$\cos(2\theta) = 1 - 2\sin^2\theta$ from Chapter~10, the equation becomes
$1 - 2\sin^2\theta + \sin\theta = 0$, in which every term now involves
sine alone: a mixed trigonometric expression has been transformed into a
quadratic equation in disguise. Substituting $u = \sin\theta$ reduces it
to $2u^2 - u - 1 = 0$, precisely the equation solved in the previous
section, so the same solution families follow immediately once $u$ is
translated back into $\sin\theta$.
\end{example}

This example captures the central purpose of the identity toolkit: an
identity is not primarily an object of memorization but a transformation
rule, valuable because it converts a difficult equation into one that can
already be solved. Much of mathematics proceeds by replacing an unfamiliar
problem with a familiar one, and the identities of Chapters~8 through 10
are the machinery that makes such replacements possible here.

\section{Extraneous Solutions and the Importance of Checking}

Every transformation used so far --- factoring, substitution, applying an
identity --- has relied on the assumption that the transformed equation has
exactly the same solutions as the original one. That assumption is often
correct, but not always. Certain algebraic operations can introduce
solutions that were not present in the original problem, called
\emph{extraneous solutions}: a value that genuinely solves a transformed
equation but fails to solve the equation from which that transformed
equation was derived. Because trigonometric equations often require
substantial manipulation before becoming solvable, this deserves careful
attention, and the safest habit is simple: every candidate solution should
be checked in the original equation before being accepted.

\begin{example}
Consider $\sin\theta = \cos\theta$. Squaring both sides gives
$\sin^2\theta = \cos^2\theta$, and using the Pythagorean identity,
$\sin^2\theta = 1 - \sin^2\theta$, so $2\sin^2\theta = 1$ and $\sin\theta =
\pm\sqrt2/2$. These values produce four candidate solutions in a full
revolution. The original equation, however, is satisfied only when
$\sin\theta$ and $\cos\theta$ share the same sign, so the angles lying in
Quadrants~II and IV, where sine and cosine have opposite signs, must be
discarded even though they satisfy the squared equation. Squaring
introduced possibilities that were not present in the original statement,
and the only reliable way to detect this is to substitute each candidate
back into the equation as originally posed.
\end{example}

Checking a candidate solution is not a tedious final formality; it is part
of the solution process itself. Some operations, such as factoring,
preserve the solution set automatically as an equation is rewritten, but
others --- squaring both sides, multiplying by an unknown quantity, or
clearing a denominator --- can enlarge it. The resulting equation may be
easier to solve, but it is no longer guaranteed to describe exactly the
original situation, and checking is what restores that guarantee. For this
reason a candidate solution is best regarded as provisional until it has
actually been verified.

\section{Solving on a Restricted Domain}

The equations solved so far sought every solution, and the answers
therefore involved an integer parameter $k \in \mathbb{Z}$, appropriate
when the domain is understood to be all real numbers. Many applications
instead specify a particular interval: an engineer may be interested only
in the first cycle of an oscillation, a navigator may need a bearing
between $0^\circ$ and $360^\circ$, or a problem may simply ask for
solutions on $[0, 2\pi)$. The underlying mathematics is unchanged; only the
final step differs, since the task becomes identifying which members of
the general solution family lie inside the specified interval, producing a
finite answer rather than an infinite one.

\begin{example}
Solve $\cos\theta = -\tfrac12$ for $0 \le \theta < 2\pi$. The reference
angle is $\pi/3$, and since cosine is negative in Quadrants~II and III, the
solutions within one revolution are $\theta = 2\pi/3$ and $\theta =
4\pi/3$. Were the general solution requested, $2\pi k$ would be appended
to each; the restricted domain instead says to stop here, so the complete
answer is
\[
\boxed{\theta = \frac{2\pi}{3},\ \frac{4\pi}{3}.}
\]
\end{example}

Restricted-domain problems are often easier to visualize geometrically than
their general-solution counterparts: rather than imagining infinitely many
revolutions around the circle, only one revolution needs to be pictured,
with the sign of the relevant trigonometric function identifying the
correct quadrants and the reference angle fixing the distance from the
nearest axis. The periodic structure remains present in the background,
but it is temporarily suppressed rather than eliminated.

\section{Applications of Trigonometric Equations}

The appearance of infinitely many solutions is not merely a mathematical
curiosity; it reflects the behavior of physical systems that repeat.
Suppose a tide reaches a particular height at noon: if the tidal cycle
repeats, that height recurs later as well, so a single solution is not
enough, and what matters is the entire sequence of times at which the
event occurs. The same reasoning governs sound waves, alternating
electrical current, planetary motion, and seasonal cycles. A model such as
$h(t) = 4 + 2\sin t$ describes a quantity oscillating continuously, and
asking when $h(t) = 5$ is asking for the times at which this graph
intersects a horizontal line --- a trigonometric equation of exactly the
kind this chapter has developed methods for. What began as geometry on the
unit circle becomes, by this point, a working tool for analyzing recurring
events throughout the natural world.

\section{Looking Back and Looking Ahead}

Part III has told a complete story. Chapter~8 began with rotations and
showed that angle addition arises from the composition of transformations,
leading to the sum and difference formulas. Chapter~9 showed that the unit
circle imposes fundamental constraints on admissible trigonometric values,
producing the Pythagorean identities and their consequences. Chapter~10
showed that a large collection of seemingly independent formulas can be
generated from that same small set of underlying principles. This chapter
has finally put those principles to work: the identities developed earlier
were not endpoints but tools, meant all along to transform equations into
forms that could be solved. Seen as a whole, Part III has been an extended
study of transformation --- rotations transformed points, identities
transformed expressions, and algebraic substitutions transformed equations
--- with the same goal recurring at every stage: replace a difficult object
with an equivalent but more useful one.

The next part of the book broadens its scope considerably. Trigonometry
has so far lived primarily inside the unit circle and the right triangle;
the chapters ahead move beyond that setting to study arbitrary triangles,
vector geometry, alternative coordinate systems, and eventually complex
numbers. The first question to ask is a simple one: what happens when the
triangle is no longer right? The answer leads to two of the most important
results in all of trigonometry, the Law of Sines and the Law of Cosines,
which extend the ideas developed throughout this book to a far wider class
of geometric problems and open Part IV.

\section{Exercises}
\begin{exercise}
Solve $\sin\theta = \dfrac{\sqrt3}{2}$. Give the general solution.
\end{exercise}

\begin{exercise}
Solve $\cos\theta = -\dfrac{\sqrt2}{2}$. Give the general solution.
\end{exercise}

\begin{exercise}
Solve $\tan\theta = 1$. Give the general solution.
\end{exercise}

\begin{exercise}
Solve $2\sin\theta - 1 = 0$. Give the general solution.
\end{exercise}

\begin{exercise}
Solve $\sin\theta(2\cos\theta + 1) = 0$. Give the general solution.
\end{exercise}

\begin{exercise}
Solve $2\sin^2\theta + \sin\theta - 1 = 0$. Give the general solution.
\end{exercise}

\begin{exercise}
Solve $\cos(2\theta) - \cos\theta = 0$, using an identity to reduce the
equation to a single function.
\end{exercise}

\begin{exercise}
Solve $\sin\theta = \dfrac12$ on the interval $0 \le \theta < 2\pi$.
\end{exercise}

\begin{exercise}
Explain why $\theta = \arcsin(a)$ is generally not the complete solution of
$\sin\theta = a$.
\end{exercise}

\begin{exercise}
A periodic signal is modeled by $V(t) = 5\sin t$. Find all times for which
$V(t) = 5/2$, expressed as a general solution set.
\end{exercise}
